% 对应单元：QM-C03-S01--QM-C03-S03
% 来源：Shankar 第 3 章第 3.1--3.3 节及学习对话中的自编检查
% 人工核验：用户作答后完成符号、传播方向、波矢几何、双缝组合规则与 Born 规则订正

\section{第 3.1 节自编检验题与订正}
\label{exercise-set:QM-C03-S01}

\exercisemeta
  {QM-C03-S01-EX}
  {2026-09-04}
  {Shankar，第 3 章 \S 3.1；本节无独立教材习题}
  {完成一组自编综合检查并逐项订正}
  {初始数据、平面波参数与波矢几何已核对}

\exercisequestion{QM-C03-S01-E01}{自编检验题：经典波的初始数据与平面波参数}

\textbf{对应知识点：}\relatedtopic{QM-C03-S01-T01}{经典粒子、经典波与初始数据}、\relatedtopic{QM-C03-S01-T02}{一维平面波、相位与传播速度}及\relatedtopic{QM-C03-S01-T03}{复数表示、强度与三维波矢}。

\textbf{来源说明。}教材第 3.1 节没有独立习题。本题由学习对话生成，用于检查本节三项核心内容，不是教材原题。

\begin{enumerate}[leftmargin=*]
  \item 说明为什么只给一维波的初始形状 $\psi(x,0)$，一般不能唯一确定它的后续运动。
  \item 对平面波
  \[
    \psi(x,t)=3\ee^{\ii(4x-12t)},
  \]
  求 $A,k,\omega,\lambda,T$、传播方向、相速度大小与强度。
  \item 若三维平面波的波矢为 $\bm k=(1,2,2)$，求 $\abs{\bm k}$、波长和等相位面的单位法向量。
\end{enumerate}

\begin{checkedsolution}
只给初始波形不能确定每一点的初始变化率。例如
\[
  \psi_+(x,t)=f(x-ct),
  \qquad
  \psi_-(x,t)=f(x+ct)
\]
在 $t=0$ 时具有相同的 $f(x)$，但初始时间导数分别为 $-cf'(x)$ 与 $cf'(x)$，所以分别向右和向左传播。还必须给出 $\dot\psi(x,0)$。

将第二问与 $A\ee^{\ii(kx-\omega t)}$ 比较，得到
\[
  A=3,
  \qquad
  k=4,
  \qquad
  \omega=12.
\]
因此
\begin{equation}
  \lambda=\frac{2\pi}{k}=\frac{\pi}{2},
  \qquad
  T=\frac{2\pi}{\omega}=\frac{\pi}{6},
  \qquad
  \abs{v_{\mathrm{ph}}}=\frac{\omega}{k}=3.
  \label{eq:QM-C03-S01-E01-plane-wave-results}
\end{equation}
固定相位满足 $4x-12t=\phi_0$，即 $x=3t+\phi_0/4$，故波向 $+x$ 传播。其强度为
\[
  I=\abs{\psi}^2=\abs{A}^2=9.
\]
指数中时间项的系数是 $-\omega$，所以 $-12t$ 对应 $\omega=12$；不能把外面的负号再次并入频率。

第三问中
\[
  \abs{\bm k}=\sqrt{1^2+2^2+2^2}=3,
  \qquad
  \lambda=\frac{2\pi}{\abs{\bm k}}=\frac{2\pi}{3}.
\]
$\bm k$ 本身就是等相位面的法向量，其单位法向量为
\begin{equation}
  \boxed{
  \widehat{\bm k}
  =\frac{\bm k}{\abs{\bm k}}
  =\frac13(1,2,2).}
  \label{eq:QM-C03-S01-E01-unit-normal}
\end{equation}
\end{checkedsolution}

\textbf{检查点：}先把指数严格写成 $kx-\omega t$ 再读取参数；传播方向由固定相位点的移动判断。三维中 $\bm k$ 是波面的法向量，而不是位于波面内的平行向量。

\section{第 3.2 节自编检验题与订正}
\label{exercise-set:QM-C03-S02}

\exercisemeta
  {QM-C03-S02-EX}
  {2026-09-05}
  {Shankar，第 3 章 \S 3.2；本节无独立教材习题}
  {完成一组自编综合检查并逐项订正}
  {振幅叠加、条纹条件、经典概率加法与相干性边界已核对}

\exercisequestion{QM-C03-S02-E01}{自编检验题：双缝干涉与经典概率加法}

\textbf{对应知识点：}\relatedtopic{QM-C03-S02-T01}{经典波的双缝干涉}、\relatedtopic{QM-C03-S02-T02}{经典粒子的概率相加}及\relatedtopic{QM-C03-S02-T03}{实验判据与相干性边界}。

\textbf{来源说明。}教材第 3.2 节没有独立习题。本题由学习对话生成，用于检查本节三项核心内容，不是教材原题。

\begin{enumerate}[leftmargin=*]
  \item 两缝在屏上产生等振幅复波
  \[
    \phi_j=A\ee^{\ii(kr_j-\omega t)}.
  \]
  从振幅叠加出发推导两缝同时打开时的强度，并明确指出干涉项。
  \item 判断 $\Delta r=2\lambda$ 与 $\Delta r=5\lambda/2$ 分别对应明纹还是暗纹，并给出强度。若其他条件不变而缝距 $d$ 加倍，条纹间距如何变化？
  \item 若经典粒子只开两缝时在某处的到达概率分别为 $P_1=0.2$、$P_2=0.3$，求两缝同时打开时的概率，并说明为什么经典模型不允许 $P_{1+2}<P_1$。最后说明“未观察到稳定干涉条纹”为何不能单独证明对象是经典粒子。
\end{enumerate}

\begin{checkedsolution}
两缝同时打开时先叠加振幅：
\[
  \phi_{1+2}=\phi_1+\phi_2.
\]
因此
\begin{align}
  I_{1+2}
  &=\abs{\phi_1+\phi_2}^2 \notag\\
  &=\abs{\phi_1}^2+\abs{\phi_2}^2
    +\underbrace{\phi_1^*\phi_2+\phi_2^*\phi_1}_{\text{干涉项}} \notag\\
  &=2\abs{A}^2\left[1+\cos(k\Delta r)\right] \notag\\
  &=4\abs{A}^2\cos^2\left(\frac{k\Delta r}{2}\right).
  \label{eq:QM-C03-S02-E01-wave-intensity}
\end{align}
干涉项不是 $\phi_1\phi_2$；复振幅取模平方时必须出现共轭。

当 $\Delta r=2\lambda$ 时，两波相位差为 $4\pi$，故为明纹，
\[
  I_{1+2}=4\abs{A}^2.
\]
当 $\Delta r=5\lambda/2$ 时，相位差为 $5\pi$，故为暗纹，
\[
  I_{1+2}=0.
\]
远屏近轴下 $\Delta x\simeq\lambda D/d$，所以把 $d$ 加倍会使条纹间距减半。

经典粒子的两条路径互斥，概率直接相加：
\[
  P_{1+2}=P_1+P_2=0.2+0.3=0.5.
\]
由于 $P_2\geq0$，必有 $P_{1+2}\geq P_1$，所以经典粒子模型不允许打开第二条路径后把该处概率降到 $P_1$ 以下。

不过，没有稳定条纹也可能是两束波的相对相位随机变化。此时干涉项在时间平均或系综平均下为零，平均强度同样表现为 $I_1+I_2$。因此“无可见条纹”不是经典粒子性的充分条件，还需控制并检验相干性。
\end{checkedsolution}

\textbf{检查点：}先分清参与相加的是振幅还是概率。波的交叉项来自 $\abs{\phi_1+\phi_2}^2$；经典互斥事件则服从概率加法。实验上还必须排除失相干造成的条纹消失。

\section{第 3.3 节自编检验题与订正}
\label{exercise-set:QM-C03-S03}

\exercisemeta
  {QM-C03-S03-EX}
  {2026-09-06}
  {Shankar，第 3 章 \S 3.3；本节无独立教材习题}
  {完成一组自编综合检查并逐项订正}
  {光子量子关系、单光子干涉、Born 规则与统计层次已核对}

\exercisequestion{QM-C03-S03-E01}{自编检验题：单光子探测、干涉与 Born 规则}

\textbf{对应知识点：}\relatedtopic{QM-C03-S03-T01}{光子能量与动量}、\relatedtopic{QM-C03-S03-T02}{单光子双缝}及\relatedtopic{QM-C03-S03-T03}{Born 规则与统计条纹}。

\textbf{来源说明。}教材第 3.3 节没有独立习题。本题由学习对话生成，用于检查本节三项核心内容，不是教材原题。

\begin{enumerate}[leftmargin=*]
  \item 一束理想单色光的频率保持不变，平均光强降为原来的 $1/5$。说明单个光子的能量和平均到达率如何变化。若随后把波长减半，单个光子的动量如何变化？
  \item 双缝在位置 $x_0$ 的概率振幅为 $\psi_1=A$、$\psi_2=-A$，在位置 $x_1$ 的概率振幅为 $\psi_1=A$、$\psi_2=\ii A$。分别计算两个位置在只开各缝和两缝同时打开时的概率密度，并说明为什么某处恰有 $\rho_{1+2}=\rho_1+\rho_2$ 不能证明没有干涉。
  \item 某一维归一化态在宽度 $\Delta x=10^{-3}\,\mathrm{m}$ 的小区间内近似满足 $\abs{\psi(x)}^2=4\,\mathrm{m}^{-1}$。求一次探测落入该区间的概率。
  \item 说明“一次只有一个光子通过装置”的实验排除了什么解释，并区分单次探测结果与大量重复后的计数分布。最后说明为什么纯经典波模型虽能给出平均条纹，却仍不足以描述本节全部现象。
\end{enumerate}

\begin{checkedsolution}
由
\[
  E_\gamma=\hbar\omega,
  \qquad
  \mathcal P=\dot N\hbar\omega
\]
可知，频率不变时单个光子的能量不变，而平均到达率随平均光强一起降为原来的 $1/5$。又因为
\[
  p_\gamma=\frac{h}{\lambda},
\]
波长减半会使单个光子的动量加倍。

在 $x_0$ 处，只开任一条缝时均有
\[
  \rho_1=\rho_2=\abs{A}^2,
\]
而两缝同时打开时
\[
  \rho_{1+2}=\abs{A-A}^2=0.
\]
在 $x_1$ 处，同样有 $\rho_1=\rho_2=\abs{A}^2$，但
\begin{align}
  \rho_{1+2}
  &=\abs{A+\ii A}^2
    =\abs{A}^2\abs{1+\ii}^2
    =2\abs{A}^2.
  \label{eq:QM-C03-S03-E01-quarter-phase}
\end{align}
此时数值上恰有 $\rho_{1+2}=\rho_1+\rho_2$，原因是
\[
  2\operatorname{Re}(A^*\ii A)=0.
\]
这只说明该位置的相位差使交叉项为零，不能说明叠加规律中不存在干涉项；换到其他位置，相位差通常不同。

第三问中概率密度在小区间内近似不变，所以
\begin{equation}
  \Pr(\text{落入该区间})
  \simeq \abs{\psi(x)}^2\Delta x
  =4\times10^{-3}
  =\boxed{\frac{1}{250}}.
  \label{eq:QM-C03-S03-E01-small-interval-probability}
\end{equation}

把入射强度降到一次只有一个光子通过，排除了“条纹由不同光子之间碰撞或相互扰动产生”的解释。单次实验只在某处记录一个完整的局域能量转移事件；大量同样制备的独立实验所形成的归一化直方图，才趋近 Born 规则给出的 $\abs{\psi(x)}^2$。

纯经典波模型能利用场振幅叠加给出平均干涉包络，但把场能量视为连续可分，不能单独说明探测器为何以能量 $\hbar\omega$ 的局域完整事件响应。量子描述用复振幅保留干涉结构，同时用 Born 规则解释离散事件的统计分布。
\end{checkedsolution}

\textbf{检查点：}固定频率时，减弱光强改变平均事件率而不改变单个光子的能量。判断干涉不能只看某一个位置是否满足概率相加，而应检查作为函数的完整分布和振幅组合规则。单次完整探测与大量重复形成条纹属于不同统计层次。
